% Source: https://artofproblemsolving.com/community/c3h1886526_geometry
% Stable identity: aops-topic-1886526-q01
Take side length $2$: $A=(0,0)$, $B=(2,0)$, $C=(2,2)$, $D=(0,2)$, $E=(1,0)$. The circle is $x^2+y^2=4$, and the points of segment $EC$ are $(1+t,2t)$ with $0\le t\le1$, where $t$ is the fraction of the way from $E$ to $C$. Substituting,
\[
(1+t)^2+4t^2=4\iff5t^2+2t-3=0\iff(5t-3)(t+1)=0,
\]
so $t=\dfrac{3}{5}$ and $F=\left(\dfrac{8}{5},\dfrac{6}{5}\right)$. Since $F$ is $\dfrac{3}{5}$ of the way from $E$ to $C$,
\[
\dfrac{EF}{FC}=\dfrac{\dfrac{3}{5}}{\dfrac{2}{5}}=\boxed{\dfrac{3}{2}}.
\]

\begin{asy}
usepackage("amsmath");
size(5cm);
defaultpen(fontsize(10pt));
pair A=(0,0),B=(2,0),C=(2,2),D=(0,2),Ee=(1,0),F=(1.6,1.2);
draw(arc(A,2,0,90),gray(0.5));
draw(A--B--C--D--cycle,linewidth(0.9));
draw(Ee--C,linewidth(0.9));
draw(A--F,blue+dashed);
dot(Ee);dot(F,red);dot(A);
label("$A$",A,SW);label("$B$",B,SE);label("$C$",C,NE);label("$D$",D,NW);label("$E$",Ee,S);label("$F$",F,dir(160),red);
label("$1$",(0.5,0),S);label("$1$",(1.5,0),S);label("$2$",(2,1),E);
label("$2$",(A+F)/2,NW,blue);
label("\mbox{}",(min(currentpicture,user=true)+max(currentpicture,user=true))/2,black); // keeps the text after the figure black when LaTeX typesets the labels inline
\end{asy}

(Check: $AF^2=\dfrac{64}{25}+\dfrac{36}{25}=4$, so $AF=2=AB$.)
